The brain has evolved to shape animal behaviour based on a lifetime of
interaction, learning and memory.  By contrast—and historically, of
necessity—much of our understanding of brain function has derived from
short-duration controlled experiments. While these experiments have revealed
key functional building blocks of brain function, they cannot explore the full
gamut of ethology; and, in any case, they do not tell us how the blocks fit
together to shape lifelong behaviour.

Recent advances in recording, video, and signal-processing technology now make
it possible to monitor animals within naturalistic, long-duration and continual
experimental settings. We have helped to pioneer this transition, by developing
the Aeon platform—a modular environment of mechatronic elements, monitoring and
recording pathways, and data curation tools.  Aeon enables a new type of
neuroscience study. In place of a tightly controlled single stimulus-behaviour
contingency, we can define multiple contingencies within a single environment
that animals explore, learn about, and engage with over weeks.

Central to Aeon is the vision of open science, with the design of the platform
and the data obtained within it to be shared widely.  This project aims to
extend that vision.  We will develop complementary platform technologies for
data analysis, visualisation, real-time feedback and sharing.  To do so, we
must overcome challenges raised by the scale and structure of such experiments.
Analyses must track changes in the statistics of underlying signals, generating
stable and robust estimates of both behavioural state and neural activity.  We
will integrate real-time machine-learning functionality into an experimental
control platform called Bonsai, developed by our industrial partners.  A
particular challenge is posed by monitoring the activity of nerve cells or
neurons.  Modern electrodes record brief electrical pulses ("spikes") from
hundreds of separate neurons, but assigning each pulse to its source neuron
("spike sorting") robustly over weeks, is made difficult by the likelihood that
the recording electrodes move within the brain tissue by many microns as well
as other factors.  Solving this challenge will be a particular focus of the
work, and one that is essential to realising the transformative potential of
long-duration experiments.  Finally, the data volumes involved (approaching
hundreds of terabytes) preclude simple sharing.  Instead we will develop a
remote-based framework to allow users to explore, visualise, and analyse
behavioural and neural datasets where they are stored; allowing the data to
uploaded once but used repeatedly.

Taken together, these tools will provide the data and analytic frameworks
needed to study  how the brain shapes exploration, learning and decisions over
the timescales that are relevant for lifelong adaptive behaviour.  Their
application will take us from a piecemeal understanding of separate brain
functions to a more wholistic view of how neural systems work and develop
together, an understanding ultimately vital to characterising developmental and
aging-related mental processes, as well as their dysfunction.

As such, this project aligns directly with the BBSRC's transformative
technologies priority within the theme of advancing the frontiers of bioscience
discovery.  Indirectly, application of these tools will advance our
understanding of the rules of life (specifically, of lifelong behaviour) and so
promote the bioscience of an integrated understanding of health.


